% ============================================================================
% CAPÍTULO 7: DISCUSIÓN
% ============================================================================
%
% REQUISITOS SEGÚN LA FACULTAD:
%
% Realiza aportaciones al conocimiento a partir de los resultados más relevantes.
% Discute los resultados considerando otras investigaciones consultadas.
% Identifica logros, riesgos y limitaciones de la investigación.
%
% ============================================================================

\chapter{Discusión}
\label{cap:discusion}

% INSTRUCCIONES GENERALES:
% - Sé CONCISO y CLARO
% - Conecta directamente con objetivos e hipótesis
% - Incluye SUSTENTO ESTADÍSTICO (valores p, coeficientes, etc.)
% - NO introduzcas información nueva
% - NO repitas la discusión completa
% - Enfócate en QUÉ se encontró, no en POR QUÉ
% - Extensión típica: 3-5 páginas

[Párrafo introductorio que contextualiza brevemente]

Con base en los resultados obtenidos y su discusión, se establecen las siguientes conclusiones:

% ----------------------------------------------------------------------------
% SECCIÓN 8.1: Conclusiones Generales
% ----------------------------------------------------------------------------
\section{Conclusiones Generales}
\label{sec:conclusiones_generales}

\begin{enumerate}
    \item Esta investigación demostró que \textit{[conclusión principal 1]} con evidencia estadística significativa (\textit{[estadístico]}, $p$ = \textit{[valor]}), lo cual confirma \textit{[implicación]}.
    
    \item Se identificó que \textit{[conclusión principal 2]}, explicando el \textit{[\%]} de la varianza ($R^2$ = \textit{[valor]}, $p < $ \textit{[valor]}), lo cual indica \textit{[implicación]}.
    
    \item Los resultados confirmaron la hipótesis de investigación al demostrar que \textit{[relación entre variables]} (\textit{[estadístico]}, $p$ = \textit{[valor]}), estableciendo que \textit{[conclusión teórica]}.
    
    \item Se validó el instrumento \textit{[nombre]} para la población de \textit{[población]}, obteniendo adecuadas propiedades psicométricas (IVC = \textit{[valor]}, $\alpha$ = \textit{[valor]}), lo cual constituye una aportación metodológica para futuras investigaciones.
    
    \item La prevalencia de \textit{[fenómeno estudiado]} en la población estudiada fue del \textit{[\%]} (IC 95\%: \textit{[inf-sup]}), lo cual \textit{[comparación con otros estudios o implicación]}.
\end{enumerate}

% ----------------------------------------------------------------------------
% SECCIÓN 8.2: Conclusiones por Objetivo Específico
% ----------------------------------------------------------------------------
\section{Conclusiones por Objetivo Específico}
\label{sec:conclusiones_por_objetivo}

% INSTRUCCIONES:
% - Presenta una conclusión para CADA objetivo específico
% - Incluye sustento estadístico
% - Sé específico sobre QUÉ se logró

\subsection{Objetivo Específico 1}
\label{subsec:conclusion_oe1}

\textbf{Objetivo:} \textit{[Enunciado del OE1]}

\textbf{Conclusión:} 

Se logró caracterizar \textit{[población/fenómeno]} encontrando que \textit{[hallazgo específico con datos]}. Los participantes presentaron una media de \textit{[variable]} de \textit{[valor]} $\pm$ \textit{[DE]}, con un rango de \textit{[mín-máx]}. La mayoría (\textit{[\%]}) se caracterizó por \textit{[característica]}.

\subsection{Objetivo Específico 2}
\label{subsec:conclusion_oe2}

\textbf{Objetivo:} \textit{[Enunciado del OE2]}

\textbf{Conclusión:}

Se determinó que \textit{[variable independiente]} se relaciona significativamente con \textit{[variable dependiente]} ($r$ = \textit{[valor]}, $p$ = \textit{[valor]}). Esta relación fue \textit{[directa/inversa]} y de magnitud \textit{[débil/moderada/fuerte]}, indicando que \textit{[interpretación concisa]}.

\subsection{Objetivo Específico 3}
\label{subsec:conclusion_oe3}

\textbf{Objetivo:} \textit{[Enunciado del OE3]}

\textbf{Conclusión:}

El análisis multivariado identificó que \textit{[predictores]} son predictores significativos de \textit{[variable dependiente]}, explicando el \textit{[\%]} de su varianza ($R^2$ ajustado = \textit{[valor]}, $p < $ \textit{[valor]}). El predictor más fuerte fue \textit{[variable]} (B = \textit{[valor]}, $p$ = \textit{[valor]}).

\subsection{Objetivo Específico 4}
\label{subsec:conclusion_oe4}

\textbf{Objetivo:} \textit{[Enunciado del OE4]}

\textbf{Conclusión:}

\textit{[Conclusión específica para este objetivo con sustento estadístico]}

[Continúa con todos tus objetivos específicos]

% ----------------------------------------------------------------------------
% SECCIÓN 8.3: Conclusión sobre la Hipótesis
% ----------------------------------------------------------------------------
\section{Conclusión sobre la Hipótesis de Investigación}
\label{sec:conclusion_hipotesis}

\textbf{Hipótesis de Investigación (Hi):}

\textit{[Enunciado completo de tu hipótesis]}

\textbf{Conclusión:}

Con base en los resultados obtenidos, se \textit{[CONFIRMA/RECHAZA]} la hipótesis de investigación. 

\textbf{Evidencia:}

\begin{itemize}
    \item Prueba estadística utilizada: \textit{[nombre de la prueba]}
    \item Valor del estadístico: \textit{[valor]}
    \item Valor de p: \textit{[valor]}
    \item Tamaño del efecto: \textit{[si aplica, reporta Cohen's d, OR, R², etc.]}
\end{itemize}

\textit{[Si se CONFIRMA:]}
Los datos proporcionan evidencia suficiente para afirmar que \textit{[relación propuesta en hipótesis]}, por lo tanto, se acepta la hipótesis de investigación y se rechaza la hipótesis nula.

\textit{[Si se RECHAZA:]}
Los datos no proporcionan evidencia suficiente para afirmar que \textit{[relación propuesta]}, por lo tanto, se rechaza la hipótesis de investigación y se retiene la hipótesis nula. Esto sugiere que \textit{[posible explicación breve]}.

% ----------------------------------------------------------------------------
% SECCIÓN 8.4: Aportaciones de la Investigación
% ----------------------------------------------------------------------------
\section{Aportaciones de la Investigación}
\label{sec:aportaciones}

Esta investigación realizó las siguientes aportaciones al campo de \textit{[tu área]}:

\subsection{Aportación Teórica}
\label{subsec:aportacion_teorica}

\begin{itemize}
    \item Se generó conocimiento nuevo sobre \textit{[aspecto específico]} en \textit{[población/contexto]}, el cual no había sido previamente documentado.
    
    \item Se aportó evidencia que \textit{[apoya/cuestiona/amplía]} la teoría de \textit{[teoría]} al demostrar que \textit{[hallazgo]}.
    
    \item Se identificaron \textit{[mecanismos/relaciones/factores]} que contribuyen a la comprensión de \textit{[fenómeno]}.
\end{itemize}

\subsection{Aportación Metodológica}
\label{subsec:aportacion_metodologica}

\begin{itemize}
    \item Se validó el instrumento \textit{[nombre]} para su uso en \textit{[población/contexto específico]} de México, facilitando futuras investigaciones en este campo.
    
    \item Se propuso una estrategia metodológica para \textit{[proceso específico]} que puede ser replicada en estudios similares.
\end{itemize}

\subsection{Aportación Práctica}
\label{subsec:aportacion_practica}

\begin{itemize}
    \item Los resultados permiten identificar \textit{[población en riesgo/factores modificables/etc.]}, lo cual puede orientar \textit{[intervenciones/políticas]}.
    
    \item Se proporcionan datos locales (Chihuahua) sobre \textit{[fenómeno]}, los cuales pueden utilizarse para \textit{[toma de decisiones/planificación/etc.]}.
    
    \item Se generó evidencia que sustenta la implementación de \textit{[intervención/programa específico]}.
\end{itemize}

% ----------------------------------------------------------------------------
% SECCIÓN 8.5: Recomendaciones
% ----------------------------------------------------------------------------
\section{Recomendaciones}
\label{sec:recomendaciones}

Con base en los hallazgos de esta investigación, se establecen las siguientes recomendaciones:

\subsection{Recomendaciones para la Práctica Profesional}
\label{subsec:recomendaciones_practica}

\begin{enumerate}
    \item Se recomienda que los profesionales de \textit{[área]} consideren \textit{[recomendación específica basada en tus hallazgos]}.
    
    \item Es importante evaluar \textit{[variable/factor]} en \textit{[población]}, dado que se identificó como predictor significativo de \textit{[resultado]}.
    
    \item Se sugiere implementar \textit{[estrategia/intervención]} para \textit{[objetivo]}, especialmente en \textit{[subgrupo identificado como de mayor riesgo]}.
\end{enumerate}

\subsection{Recomendaciones para Políticas de Salud Pública (si aplica)}
\label{subsec:recomendaciones_politicas}

\begin{enumerate}
    \item Se recomienda a las autoridades de salud de \textit{[nivel: local/estatal/nacional]} considerar \textit{[recomendación de política]} basándose en que \textit{[hallazgo que la sustenta]}.
    
    \item Sería conveniente desarrollar programas dirigidos a \textit{[población específica]} enfocados en \textit{[objetivo]}.
\end{enumerate}

\subsection{Recomendaciones para Futuras Investigaciones}
\label{subsec:recomendaciones_investigacion}

\begin{enumerate}
    \item Se recomienda replicar este estudio con un diseño \textit{[longitudinal/experimental]} para establecer relaciones causales entre \textit{[variables]}.
    
    \item Futuras investigaciones deberían incluir una muestra más amplia y diversa, incorporando \textit{[grupos no representados en este estudio]}.
    
    \item Sería valioso explorar \textit{[aspecto no abordado]} mediante \textit{[metodología sugerida]}.
    
    \item Se sugiere investigar el efecto de \textit{[intervención propuesta]} en \textit{[variable de resultado]} a través de un ensayo clínico aleatorizado.
    
    \item Es necesario realizar estudios que incluyan mediciones objetivas de \textit{[variable]}, complementando los datos de autorreporte.
    
    \item Se recomienda investigar los mecanismos mediante los cuales \textit{[variable X]} influye en \textit{[variable Y]}, posiblemente a través de \textit{[mediadores sugeridos]}.
\end{enumerate}

% ----------------------------------------------------------------------------
% SECCIÓN 8.6: Cierre
% ----------------------------------------------------------------------------
\section{Reflexión Final}
\label{sec:reflexion_final}

Esta investigación logró cumplir con los objetivos planteados y aportar conocimiento significativo sobre \textit{[tema central]}. Los hallazgos demuestran que \textit{[mensaje central de la tesis]}, lo cual tiene implicaciones importantes para \textit{[campo/práctica/teoría]}.

A pesar de las limitaciones identificadas, los resultados proporcionan una base sólida para \textit{[aplicación práctica/futuras investigaciones/desarrollo teórico]}. Se espera que este trabajo contribuya a \textit{[impacto esperado]} y motive futuras investigaciones en esta área.

El cumplimiento de los objetivos y la confirmación de la hipótesis mediante evidencia estadística robusta validan la relevancia y pertinencia de este estudio. Las aportaciones teóricas, metodológicas y prácticas realizadas representan una contribución significativa al campo de \textit{[tu área de estudio]}.

% ============================================================================
% NOTAS FINALES:
% ============================================================================
%
% RECUERDA:
% 1. Sé CONCISO y CLARO
% 2. Incluye SUSTENTO ESTADÍSTICO en cada conclusión
% 3. Conecta directamente con OBJETIVOS e HIPÓTESIS
% 4. NO introduzcas información NUEVA
% 5. Enfócate en QUÉ se encontró
% 6. Incluye aportaciones CONCRETAS
% 7. Recomendaciones deben ser ESPECÍFICAS y VIABLES
%
% ============================================================================

